\chapter{Abstract}

The objective of the thesis is to characterize three aspects of human settlements in Mexico: their geographic extent, their size distribution, and the structure of their street networks. I develop and apply mathematical methods, specifically computational geometry, graph and network theory, probability, and statistics. In the first chapter, I propose an algorithm for delimiting human settlements based on the intersection graph of the convex hull of polygons representing built-up areas. I indicate that the limits generated by the algorithm better reflect the geographic extent of cities than their administrative or economic boundaries. In the second chapter, I review the state-of-the-art statistical methods for verifying if an empirical sample follows a power-law distribution. I apply these methods to characterize the size distribution of cities in Mexico, the United States, Canada, Chile, Colombia, and Uruguay. I find that the size distribution of cities in Mexico has followed a power law over the past 20 years, extending to small urban concentrations. I also identify that the size distribution of cities in the three North American countries follows a power law, both individually and collectively. In the third chapter, I theoretically and computationally address some properties of $\beta$-skeletons. I prove that tessellations of cyclic polygons are $\beta$-skeletons. Furthermore, I show that $\beta$-skeletons approximate regular tessellations with moderately perturbed vertices well. I apply two theorems from the literature to calculate the degree and average length of $\beta$-skeletons of random points. In the final chapter, I analyze 11,091 street networks using \textit{topological} and geometric indices proposed in the literature and the degree of approximation of $\beta$-skeletons. A conclusion is that the level of approximation of $\beta$-skeletons can be used as an alternative to evaluate the similarity of a street network to a rectangular grid.